\section{¿Qué resuelve SAE?}

La Estimación en Áreas Pequeñas (Small Area Estimation, SAE) surge cuando se desea estimar un parámetro de interés para dominios muy desagregados, por ejemplo: Distritos, minicipios, etc.
\\
El problema es que en esos dominios, el tamaño muestral suele ser pequeño. Como consecuencia, los estimadores directos basados únicamente en la encuesta presentan alta variabilidad y baja precisión, incluso pueden no existir para algunas áreas no muestreadas.
\\
Formalmente, si $\theta_i$ representa el parámetro de interés del área $i$, y $\hat{\theta}_i^{\,D}$ su estimador directo, entonces el problema típico en SAE es que
\[
\operatorname{Var}\!\left(\hat{\theta}_i^{\,D}\right)
\]
puede ser grande cuando el tamaño muestral del área $i$ es pequeño.

\section{Formulación del problema}

Sea $i=1,\dots,m$ el índice de áreas pequeñas. Para cada área $i$, queremos estimar un parámetro $\theta_i$, por ejemplo una media, una proporción, una tasa o un índice de pobreza.

Disponemos de un estimador directo $y_i$, que aproxima a $\theta_i$. En SAE suele escribirse
\[
y_i = \hat{\theta}_i^{\,D}.
\]

Este estimador directo satisface, de manera aproximada,
\[
E(y_i \mid \theta_i) \approx \theta_i,
\qquad
\operatorname{Var}(y_i \mid \theta_i) = D_i,
\]


El gran problema es que cuando $D_i$ es grande, el estimador directo es inestable y poco confiable.

\section{SAE}
 SAE suele dividirse en dos grandes familias.

\subsection{Modelos a nivel unidad}

Trabajan con microdatos, es decir, con observaciones individuales $y_{ij}$.

\subsection{Modelos a nivel área}

Trabajan con un resumen por área:
\begin{itemize}
    \item un estimador directo $y_i$,
    \item su varianza $D_i$,
    \item covariables agregadas $x_i$.
\end{itemize}

El modelo clásico de Fay--Herriot pertenece a esta segunda familia.


